\section{Introduction}
\label{sec:introduction}

The fixed-shift native-operator route separates the missing-high Gram
defect into a same-residual part \(R\) and an erasure-created part \(E\).
TPC--68 identifies the signed operator certificates attached to
\(R+E\), and TPC--69 proves that the first mixed trace is a cross-key
rectangle rather than a one-key correction
\citep{WangTPC68,WangTPC69}.  At its shortest, the relevant term has two
source rows and two complete native keys:
\[
 \begin{matrix}
     &m&n\\
 \text{erasure key}&\bullet&\bullet\\
 \text{same-residual key}&\bullet&\bullet .
 \end{matrix}
\]
The erasure edge carries the residual sign; the return edge has equal
residual labels and contributes no variable target sign after the closed
product is formed.

The immediate question is arithmetic rather than matrix-theoretic:
\emph{what coordinates does an actual rectangle possess, and along which
coordinate can its M\"obius sign vary?}  The answer is more rigid than a
generic two-time correlation.  Equal residual labels at the return key
compress the two cofactors to a single divisor \(t\) of
\(h_0(n-m)\).  At the erasure key, the common cofactor content is another
divisor \(q\) of the same determinant.  After removing this content, the
only variable arithmetic sign is \(\mu(u)\mu(v)\) on a primitive
two-dimensional cofactor plane.

\subsection{Main results}

Let
\[
 A_m=mj_R+h_0=P_mR_m,\qquad
 A_n=nj_R+h_0=P_nR_n
\]
be the two targets at the same-residual key, and let
\[
 B_m=mj_E+h_0=Q_mS_m,\qquad
 B_n=nj_E+h_0=Q_nS_n
\]
be the targets at the erasure key.  The letters \(P_m,P_n,Q_m,Q_n\)
denote the actual largest-prime labels of the four physical atoms.  The
cofactors are \(R_m,R_n,S_m,S_n\).

The paper proves the following exact chain.
\begin{enumerate}[label=\textup{(\arabic*)},leftmargin=3em]
 \item The four targets obey two horizontal determinant equations, two
       vertical difference equations, and one cross determinant:
       \[
       nA_m-mA_n=nB_m-mB_n=h_0(n-m),
       \qquad A_mB_n-B_mA_n=h_0(n-m)(j_E-j_R).
       \]

 \item If \(d_m=ga\), \(d_n=gb\), \((a,b)=1\), then the same-residual
       equality \(d_mR_m=d_nR_n\) is equivalent on the squarefree carrier
       to
       \[
       R_m=bt,\qquad R_n=at.
       \]
       The horizontal determinant forces \(t\mid h_0(n-m)\).  Every
       fixed-\(t\) physical anchor lies on one affine prime ray of spacing
       \(abt\).

 \item Writing
       \[
       q=(S_m,S_n),\qquad S_m=qu,\quad S_n=qv,\quad (u,v)=1,
       \]
       the erasure determinant forces \(q\mid h_0(n-m)\).  The erased
       residual labels are unequal precisely when
       \[
       (u,v)\ne(b,a).
       \]
       Every fixed-\((q,u,v)\) physical erasure lies on one affine prime
       ray of spacing \(quv\).

 \item The exact rectangle sign is
       \[
       \mu(a)\mu(b)\mu(u)\mu(v).
       \]
       The common factors \(g,q\) disappear, and the sign is constant on
       each time ray.  Therefore an exact grouping of the literal
       rectangle sum yields a primitive two-dimensional M\"obius
       transform in \((u,v)\).
\end{enumerate}

The compression is real: \(t\) and \(q\) run only over divisors of a fixed
determinant for a fixed row pair.  It is not cancellation.  The
\(z\)-fibers still carry two affine primality conditions, a largest-prime
condition, interval restrictions, missing and first-failure masks,
complete native labels, and literal complex weights.  More importantly,
their M\"obius sign is frozen.  A Brun or Selberg upper-bound sieve can
bound the number of prime points on one ray, but it cannot make
\(\mu(u)\mu(v)\) oscillate along that ray.

\subsection{Why this is a useful boundary}

The result rules out an attractive but incorrect strategy: reparametrize
by time and hope that the target M\"obius signs fluctuate along each
prime progression.  They do not.  Any future signed gain must come from
one of two sources:
\begin{enumerate}[label=\textup{(\roman*)},leftmargin=3em]
 \item cancellation across primitive cofactor pairs \((u,v)\); or
 \item literal complex-phase cancellation already present in the physical
       kernel.
\end{enumerate}
This locates the next analytic door without claiming that the door is
open.  A natural continuation is a two-dimensional summation-by-parts
criterion for the actual ragged kernel, including its boundary and
variation costs.

\subsection{Claim discipline}

We use the following labels throughout.
\begin{itemize}
 \item \textbf{L0} denotes finite exact arithmetic and algebraic
       identities.
 \item \textbf{L1} denotes their exact identification with the complete
       fixed-\(h_0\) native carrier, with no physical label or coefficient
       removed.
 \item \textbf{L2} denotes a genuinely new cancellation estimate for the
       literal growing arithmetic family.
\end{itemize}
This paper proves L0 identities and an L1 reindexing.  It proves no L2
estimate.  In particular, divisor compression, a ray census, or an
unsigned prime upper bound is not fixed-\(h_0\) parity progress.

\subsection{Organization}

\Cref{sec:physical-interface} fixes the literal two-key carrier.
\Cref{sec:four-corner} proves the four-corner identities.
\Cref{sec:anchor-compression} compresses the same-residual anchor, and
\cref{sec:erasure-coordinates} constructs the primitive erasure
coordinates.  The ray and sign theorems are proved in
\cref{sec:rays-and-signs}.  \Cref{sec:cofactor-plane} gives the exact
M\"obius-plane transform.  \Cref{sec:sieve-boundary} records the unsigned
sieve boundary, and \cref{sec:counterexamples} gives sharp examples.
The existing-theorem mismatch, operator ledger, and complete claim
boundary occupy the final three sections.
